\documentclass[11pt,a4paper]{article}
\usepackage[utf8]{inputenc}
\usepackage[T1]{fontenc}
\usepackage[spanish]{babel}
\usepackage{hyperref}
\usepackage{geometry}
\geometry{margin=2.5cm}
\usepackage{parskip}

\title{098 la friccion quitada con firma es una decision; sin firma es un accidente}
\author{Alejandro Valenzuela Rivera\\\small Universidad Autónoma de Nuevo León}
\date{\today}
\begin{document}
\maketitle

\begin{quote}
\textbf{Sesgo declarado.} Este trabajo se sostiene desde el sesgo declarado de su autor, y no finge neutralidad. Lo que afirma se puede tumbar, y ese es el punto: el falsificador está escrito, no escondido.
\end{quote}

\section*{Resumen}
El proceso no impide quitar la friccion: la marca. Quien la quita firma con su nombre, dice que quito, que acepta y deja el hash de lo que produjo sin ella. Asi el costo deja de ser de los demas y pasa a ser suyo.

\section{Por qué se sostiene}
No se limita: se estandariza la consecuencia. Y el lugar de la firma es libre a proposito, porque si la firma depende de una herramienta, la herramienta se queda con la fama de ser la unica que puede usar la tecnologia. Un hash se verifica en cualquier parte, hasta en papel.

\section{De dónde viene}
La autoridad, y el proceso FIRMAR.prs.

\section{El caso}
Ninguna: es la regla de la firma.

\section{Método}
Una friccion quitada, con su hash, y con quien la quito.

\section{Límites, y qué lo tumba}
Una friccion quitada sin firma: nadie decidio quitarla.

\section{Lo que queda abierto}
\textquestiondown{}Que pasa si alguien quita la friccion y no firma?

\section*{Cómo se leyó esta tesis}
Esto se escribió en el repositorio man-made-empathy, del que este candado es una pieza: candados/04 el lenguaje/098 la friccion quitada con firma es una decision; sin firma es un accidente.md. La máquina que lo compuso lo hizo desde este mismo archivo, y no agregó nada que no estuviera escrito. Si algo de esto es falso, se dice en el mismo sitio, y se corrige con fecha.

\end{document}
